% Edited conversion of source p. 47.
\ReportHeading
  {Розв’язання задачі лінійної в’язкопружності для анізотропних плит з еліптичними неоднорідностями}
  {А.~О. Кошкін\textsuperscript{1}, О.~О. Стрельнікова\textsuperscript{1,2}}
  {\textsuperscript{1}Харківський національний університет радіоелектроніки; \textsuperscript{2}Інститут енергетичних машин і систем ім. А.~М. Підгорного НАН України}
  {}

Дослідження часової зміни напруженого стану багатозв’язних в’язкопружних плит після прикладання навантаження є актуальним для механіки композитів. Конструкційні елементи з локальними неоднорідностями --- включеннями, отворами та армувальними елементами --- широко використовуються в авіабудуванні, космічній техніці та будівництві, що потребує ефективних методів довготривалого розрахунку.

Розглянуто багатозв’язну анізотропну в’язкопружну плиту, в отвори якої вставлено еліптичні включення з інших матеріалів. Як окремий випадок зовнішній контур віддаляється на нескінченність, і задача переходить до нескінченної плити.

Для визначення характеристик напруженого стану у часі використано принцип Вольтерри~\cite{viscoplate:volterra1913}, за яким пружні сталі у розв’язку відповідної пружної задачі замінюються часовими операторами. Оскільки для багатозв’язної області явний розв’язок через пружні сталі недоступний, застосовано метод малого параметра. Малим параметром $\lambda$ обрано зміну коефіцієнта Пуассона у часі.

Розв’язок побудовано за допомогою комплексних потенціалів Лехницького, поданих як наближення $j$-го порядку. Для плити потенціали розкладено у ряди Лорана, а для включень --- у ряди за многочленами Фабера. Невідомі сталі визначено узагальненим методом найменших квадратів із граничних умов~\cite{viscoplate:koshkin2025}.

Після побудови пружного розв’язку здійснено перехід до в’язкопружної постановки: степені $\lambda^j$ замінено виразами, що містять реологічні сталі матеріалу плити, миттєве пружне значення коефіцієнта Пуассона та оператори Работнова~\cite{viscoplate:rabotnov1969}. Це дає змогу обчислювати згинальні моменти й поперечні сили у довільний момент часу $t$.

Чисельні дослідження проведено для алюмінію та епоксидного матеріалу. Найінтенсивніший перерозподіл моментів відбувається на початковому етапі --- протягом перших 50--100 годин для розглянутих матеріалів, після чого система прямує до стаціонарного стану. Концентрація моментів істотно залежить від відносної жорсткості включення; найбільші значення виникають для граничних випадків абсолютно жорстких і абсолютно м’яких включень, тобто отворів.

Запропонований чисельно-аналітичний підхід поєднує апарат комплексних потенціалів із методом малого параметра та може бути використаний для прогнозування довготривалої міцності анізотропних конструкцій із неоднорідностями.

\begin{thebibliography}{9}
\bibitem{viscoplate:volterra1913}
V.~Volterra, \emph{Leçons sur les fonctions de lignes}.
Paris: Gauthier-Villars, 1913, 230~p.

\bibitem{viscoplate:koshkin2025}
A.~Koshkin and O.~Strelnikova,
``Bending analysis of multiply-connected anisotropic plates with elastic inclusions,''
\emph{Bulletin of V.~N. Karazin Kharkiv National University, Series Mathematical Modeling, Information Technology, Automated Control Systems}, vol.~68, pp.~43--52, 2025.
\url{https://doi.org/10.26565/2304-6201-2025-68-04}.

\bibitem{viscoplate:rabotnov1969}
Yu.~N. Rabotnov, \emph{Creep Problems in Structural Members}.
Amsterdam: North-Holland, 1969, 822~p.
\end{thebibliography}
